\clearpage\fancyhead[L]{\small VESTA.}\hypertarget{native-366-lower}{}\label{native-366-lower}
\begin{center}\textbf{VESTA.}\par\vspace{4pt}\hrule\vspace{4pt}\fontsize{10}{12}\selectfont Göttingische gelehrte Anzeigen. Stück 205. Seite 2041\;..\;2045. 24. Dezember 1812.\par\vspace{4pt}\hrule\vspace{4pt}\end{center}

\par\indent Im Oktober dieses Jahres kam Vesta, der jüngste der neuen Planeten, zum vierten Mal seit ihrer Entdeckung in Opposition zur Sonne. Für die Theorie eines neuen Planeten ist gerade die vierte Opposition stets besonders wichtig, weil man mit ihr die Bahn erstmals allein aus Oppositionen bestimmen kann. Deshalb wurde Vesta an der hiesigen Sternwarte um die Zeit der Opposition so oft im Meridian beobachtet, wie es das ungünstige Wetter zuließ. Professor \textsc{Gauss} achtete dabei besonders auf die Deklinationen. Die \SourceLeaf{367}Rektaszensionen wurden daher nur nebenbei mitgemessen, weil man hoffen konnte, von anderen Sternwarten bessere Werte zu erhalten, als sie sich am Mauerquadranten messen lassen. Diese Beschränkung auf das, was sich mit dem Mauerquadranten gut messen lässt, war umso notwendiger, als der Planet bei Weitem weniger als halb so hell war wie bei der Opposition des Vorjahres und sich diesmal mit dem lichtschwachen Instrument oft nur mühsam beobachten ließ. Vesta wurde insgesamt fünfmal beobachtet. Die Beobachtungen vom 24. und 27. Oktober waren jedoch bei sehr ungünstigen Luftverhältnissen eher Schätzungen. Die drei Deklinationen vom 25. Oktober, 28. Oktober und 1. November sind dagegen als sehr gut anzusehen. Am 1. November wurde die Rektaszension nur nebenbei an einem einzigen Faden gemessen.

\medskip
\Needspace{12\baselineskip}\noindent\emph{Beobachtung der Vesta in Göttingen.}

\begin{center}
\fontsize{10}{12}\selectfont\setlength{\tabcolsep}{4pt}\renewcommand{\arraystretch}{1.15}\begin{tabular}{|c|c|c|c|}
\hline
1812 & Mittlere Zeit & Rektaszension & Nördliche Dekl. \\
\hline
Okt. 24 & $12^{\mathrm u}\;3^{\mathrm m}\;39^{\mathrm s}$ & . . . . & $1^\circ\,54'\,49''$ \\
25 & $11\;58\;43$ & $33^\circ\,52'\,33''1$ & $1\;50\;15.8$ \\
27 & $11\;48\;53$ & . . . . & $1\;41\;56$ \\
28 & $11\;43\;58$ & $33\;\;8\;\;1.8$ & $1\;38\;40.9$ \\
Nov. 1 & $11\;24\;21$ & $32\;\;9\;35.9$ & $1\;25\;\;0.9$ \\
\hline
\end{tabular}
\end{center}

\Needspace{13\baselineskip}\par\indent An der Seeberger Sternwarte konnte der Planet nur zweimal beobachtet werden. Herr von \textsc{Lindenau} teilte Professor \textsc{Gauss} folgende Beobachtungen mit:

\begin{center}
\fontsize{10}{12}\selectfont\setlength{\tabcolsep}{4pt}\renewcommand{\arraystretch}{1.15}\begin{tabular}{|c|c|c|c|}
\hline
1812 & Mittlere Zeit & Rektaszension & Nördliche Dekl. \\
\hline
Okt. 25 & $11^{\mathrm u}\;58^{\mathrm m}\;46^{\mathrm s}8$ & $33^\circ\,52'\,31''7$ & $1^\circ\,50'\,28''1$ \\
Nov. 1 & $11\;24\;23.1$ & $32\;\;9\;11.0$ & $1\;25\;28{,}9$ \\
\hline
\end{tabular}
\end{center}

\Needspace{17\baselineskip}\par\indent Schließlich machte Herr von \textsc{Zach} an seiner Sternwarte bei Marseille folgende Meridianbeobachtungen:

\begin{center}
\fontsize{10}{12}\selectfont\setlength{\tabcolsep}{4pt}\renewcommand{\arraystretch}{1.15}\begin{tabular}{|c|c|c|c|}
\hline
1812 & Mittlere Zeit & Rektaszension & Nördliche Dekl. \\
\hline
Okt. 21 & $12^{\mathrm u}\;18^{\mathrm m}\;21^{\mathrm s}1$ & $34^\circ\,51'\,13''2$ & $2^\circ\,\;7'\,57''5$ \\
23 & $12\;\;8\;32.3$ & $34\;21\;53.4$ & $1\;54\;39{,}9$ \\
24 & $12\;\;3\;37.5$ & $34\;\;7\;\;8.4$ & $1\;54\;25.3$ \\
26 & $11\;53\;47.5$ & $33\;37\;31.2$ & $1\;46\;15.7$ \\
28 & $11\;43\;57.4$ & $33\;\;7\;52.2$ & $1\;38\;54.5$ \\
29 & $11\;39\;\;2.4$ & $32\;53\;\;2.4$ & $1\;35\;\;0.1$ \\
\hline
\end{tabular}
\end{center}

\par\indent Beim Vergleich all dieser Beobachtungen mit den neuesten Elementen, die in der Monatl. Corresp., Band XXIV, S. 502, mitgeteilt wurden, erhöhte man die Epoche um $12'\,47''3$, um eine bessere Übereinstimmung zu erhalten. Dadurch ergaben sich folgende Unterschiede.

\medskip
\Needspace{14\baselineskip}\noindent\emph{Beobachtungen in Göttingen.} 

\begin{center}
Unterschied

\fontsize{10}{12}\selectfont\setlength{\tabcolsep}{7pt}\renewcommand{\arraystretch}{1.15}\begin{tabular}{|c|c|c|}
\hline
 & Rektaszension & Deklination \\
\hline
Okt. 24 & . . . & $-38''68$: \\
25 & $+3''04$ & $-13.26$ \\
27 & . . . & $+16.10$: \\
28 & $+6.37$ & $-13.32$ \\
Nov. 1 & $-15.16$ & $-10.17$ \\
\hline
\end{tabular}
\end{center}

\Needspace{12\baselineskip}\SourceLeaf{368}\noindent\emph{Beobachtungen auf der Seeberger Sternwarte.}

\begin{center}
Unterschied.

\fontsize{10}{12}\selectfont\setlength{\tabcolsep}{7pt}\renewcommand{\arraystretch}{1.15}\begin{tabular}{|c|c|c|}
\hline
 & Rektaszension & Deklination \\
\hline
Okt. 25 & $+6''44$ & $-23''41$ \\
Nov. 1 & $+11.71$ & $-37.70$ \\
\hline
\end{tabular}
\end{center}

\medskip
\Needspace{15\baselineskip}\noindent\emph{Beobachtungen zu Lacapellete bei Marseille.}

\begin{center}
Unterschied.

\fontsize{10}{12}\selectfont\setlength{\tabcolsep}{7pt}\renewcommand{\arraystretch}{1.15}\begin{tabular}{|c|c|c|}
\hline
 & Rektaszension & Deklination \\
\hline
Okt. 21 & $+2''73$ & $-52''74$ \\
23 & $+3.93$ & $+220.85$ \\
24 & $+3.80$ & $-17.41$ \\
26 & $+5.06$ & $-13.51$ \\
28 & $+4.80$ & $-29.19$ \\
29 & $+7.20$ & $-11.60$ \\
\hline
\end{tabular}
\end{center}

\par\indent Man sieht, dass die Rektaszensionen gut übereinstimmen. Bei den Deklinationen ist die Übereinstimmung weniger gut. Professor \textsc{Gauss} meinte, sich hier vor allem an die drei hiesigen Deklinationsmessungen halten zu müssen, die untereinander sehr gut übereinstimmen. Bei der Beobachtung von Herrn von \textsc{Zach} am 23. Oktober scheint ein Fehler von 4 Minuten unterlaufen zu sein. Unter dieser Annahme stimmen wenigstens die vier Beobachtungen vom 23., 24., 26. und 29. Oktober auch untereinander recht gut und mit den hiesigen annähernd überein. So ergab sich schließlich für die Opposition folgendes Ergebnis:

\medskip
\begin{center}\fontsize{10}{12}\selectfont\setlength{\tabcolsep}{2pt}\renewcommand{\arraystretch}{1.15}\begin{tabular}{l}
1812 Oktober $8^{\mathrm u}\;54^{\mathrm m}\;44^{\mathrm s}$\\Mittlere Zeit in Göttingen
\end{tabular}\end{center}
\begin{center}\fontsize{10}{12}\selectfont\setlength{\tabcolsep}{2pt}\renewcommand{\arraystretch}{1.15}\begin{tabular}{lcr}
Wahre Länge & $.\;.\;.\;.\;.\;.\;.\;.\;.\;.\;.$ & $32^\circ\;17'\;40''8$\\
Geozentrische Breite & $.\;.\;.\;.\;.\;.\;.\;.\;.$ & $11\;5\;32.3\;\text{südl.}$\\
\end{tabular}\end{center}

\Needspace{22\baselineskip}\par\indent Die Berichtigung der Elemente anhand der vier bisher beobachteten Oppositionen übertrug Professor \textsc{Gauss} daraufhin Herrn \textsc{Enke}, der sich bei uns mit ausgezeichnetem Erfolg dem Studium der Astronomie widmet und in astronomischen Rechnungen bereits sehr geschickt ist. Sein Ergebnis lautet:

\medskip
\begin{center}\fontsize{10}{12}\selectfont\setlength{\tabcolsep}{2pt}\renewcommand{\arraystretch}{1.15}\begin{tabular}{l}
Epoche der Länge 1814 Februar 13.\\$0^{\mathrm u}$ in Göttingen $\; .\; 154^\circ\;48'\;29''4$
\end{tabular}\end{center}
\begin{center}\fontsize{9}{11}\selectfont\setlength{\tabcolsep}{2pt}\renewcommand{\arraystretch}{1.15}\begin{tabular}{lcr}
Tägliche mittlere tropische Bewegung & $.\;.\;.\;.\;.\;.$ & $978''1642$\\
Länge des Perihels & $.\;.\;.\;.\;.\;.\;.\;.\;.\;.\;.\;.$ & $249^\circ\;38'\;31.4$\\
Länge des aufsteigenden Knotens & $.\;.\;.\;.\;.\;.\;.\;.$ & $103\;12\;25.3$\\
Neigung der Bahn & $.\;.\;.\;.\;.\;.\;.\;.\;.\;.\;.\;.\;.$ & $7\;8\;5.0$\\
Exzentrizität & $.\;.\;.\;.\;.\;.\;.\;.\;.$ & $=0.0894779=\sin 5\;8\;0.86$\\
Logarithmus der großen Halbachse & $.\;.\;.\;.\;.\;.\;.$ & $0.3731047$\\
\end{tabular}\end{center}

\Needspace{12\baselineskip}\par\indent Die nächste Opposition fällt nach der auf diesen Elementen beruhenden Rechnung von Herrn \textsc{Enke} auf den 13. Februar 1814, $10^{\mathrm u}\;21^{\mathrm m}\;26^{\mathrm s}$, bei $144^\circ\,37'\,51''$ Länge und $8^\circ\,2'\,19''$ nördlicher Breite. Die Helligkeit des Planeten wird dann $= 0.08698$ sein. Bei den ersten vier Oppositionen betrug sie

\medskip
\begin{center}\begin{tabular}{ll}$1808\;.\;.\;0.09516,$ & $1810\;.\;.\;0.06440,$\\$1811\;.\;.\;0.16540,$ & $1812\;.\;.\;0.06641$\end{tabular}\end{center}

\medskip
\noindent Eine vollständige Ephemeride für 1813 und 1814, berechnet von Herrn \textsc{Enke}, wird an anderer Stelle erscheinen.

\Needspace{39\baselineskip}\SourceLeaf{369}\par\indent Wir nutzen diese Gelegenheit, um kurz über eine wichtige Arbeit zu berichten, die Herr \textsc{Nicolai} soeben über den zweiten Kometen des Jahres 1811 abgeschlossen hat. Von diesem Kometen war in diesen Blättern schon zweimal die Rede (1811, Stück 201, und 1812, Stück 15 [S. 342 und 345 dieses Bandes]). Im letzteren Heft wurden die Ergebnisse der Rechnungen von Herrn \textsc{Nicolai} mitgeteilt. Die ausgezeichnete Genauigkeit, die die Beobachtungen dieses Kometen zuließen, und der keineswegs kurze Zeitraum, den sie abdeckten, schienen Professor \textsc{Gauss} auszureichen, um etwas über die Elliptizität der Bahn entscheiden zu können. Er ermunterte Herrn \textsc{Nicolai} daher zu dieser interessanten Untersuchung, und das Ergebnis hat dieses Urteil vollständig bestätigt. Unter allen Kometen, deren Umlaufzeit aus einer einzigen Erscheinung berechnet wurde, gehört dieser gewiss zu denen mit der am besten bestimmten Umlaufzeit. Das Ergebnis von Herrn \textsc{Nicolai} lautet:

\Needspace{22\baselineskip}\medskip
\begin{center}\fontsize{10}{12}\selectfont\setlength{\tabcolsep}{2pt}\renewcommand{\arraystretch}{1.15}\begin{tabular}{l}
Durchgang durch das Perihel 1811 Nov. 11.\\$0^{\mathrm u}\;26^{\mathrm m}\;4^{\mathrm s}$
\end{tabular}\end{center}
\begin{center}\fontsize{9}{11}\selectfont\setlength{\tabcolsep}{2pt}\renewcommand{\arraystretch}{1.15}\begin{tabular}{lcr}
Länge des Perihels & $.\;.\;.\;.\;.\;.\;.\;.\;.\;.\;.\;.$ & $47^\circ\;27'\;27''1$\\
Länge des aufsteigenden Knotens & $.\;.\;.\;.\;.\;.\;.\;.\;.$ & $93\;1\;52.1$\\
\end{tabular}\end{center}
\noindent beide vom mittleren Äquinoktium des 1. Januar 1812 aus gezählt
\begin{center}\fontsize{9}{11}\selectfont\setlength{\tabcolsep}{2pt}\renewcommand{\arraystretch}{1.15}\begin{tabular}{lcr}
Neigung der Bahn & $.\;.\;.\;.\;.\;.\;.\;.\;.\;.\;.\;.\;.\;.$ & $31^\circ\;17'\;11''0$\\
Logarithmus des kleinsten Abstands & $.\;.\;.\;.\;.\;.\;.$ & $0.1992359$\\
Exzentrizität & $.\;.\;.\;.\;.\;.\;.\;.\;.\;.\;.\;.\;.\;.\;.\;.\;.$ & $0.98271088$\\
Große Halbachse & $.\;.\;.\;.\;.\;.\;.\;.\;.\;.\;.\;.\;.\;.$ & $91.5088$\\
Siderische Umlaufzeit & $.\;.\;.\;.\;.\;.\;.\;.\;.\;.\;.\;.$ & $875.4\;\text{Jahre.}$\\
\end{tabular}\end{center}

\par\indent Eine ausführliche Darstellung dieser Arbeit behält sich Herr \textsc{Nicolai} für eine andere Veröffentlichung vor. Sie zeigt die ausgezeichnete Übereinstimmung dieser Ellipse mit den Beobachtungen und die Unmöglichkeit, diese mit einer parabolischen Bahn zu vereinbaren.
